\begin{enumerate}
    \item $V, W$をｔ複素線形空間とし、$V, W$の零ベクトルをそれぞれ$\bm{0}_V, \bm{0}_W$とする。$V$から$W$への線形写像$T\colon V\to W$に対して、$T$の像を$\Im{(T)}$とし、$T$の核を$\Ker{(T)}$とする。$\Im{(T)}$および$\Ker{(T)}$の定義をそれぞれ述べなさい。
    \item $\C$を複素数体とする。線形写像$S\colon \C^3\to \C^4$および$T\colon \C2^\to \C^4$を以下で定める。
\begin{equation*}
    S\left(\begin{pmatrix}x_1& x_2& x_3\end{pmatrix}\right)= 
    \begin{pmatrix}
        3x_1- x_2+ 5x_3\\
        x_1+ x_2+ 3x_3\\
        x_1+ 2x_3\\
        x_1- 3x_2- x_3
    \end{pmatrix},\ 
    T\left(\begin{pmatrix}x_1\\ x_2\end{pmatrix}\right)= 
    \begin{pmatrix}
        x_1- 3x_2\\
        3x_1+ x_2\\
        x_1- 2x_2\\
        -5x_1- x_2
    \end{pmatrix}
\end{equation*}
\begin{enumerate_alpha}
    \item $\Ker{(S)}= \begin{Bmatrix}\alpha\begin{pmatrix}-2\\ -1\\ 1\end{pmatrix}; \alpha\in \C\end{Bmatrix}$を示しなさい。
    \item $\Im{(S)}$と$\Im{(S)}$の和空間を$W= \Im{(S)}+ \Ker{(T)}$とおく。$W$の次元を求めなさい。またその根拠も述べなさい。

\end{enumerate}